\section{Frobenius at a conductor prime}\label{sec:self-conductor}
The toric row $cz(F)-bz(T)$ needs $b=\beta(F)$ even where beta ramifies. An unramified test prime formula is insufficient. The final calculation includes the self conductor unit and converts the old twisted inertia normalization to the ordinary Artin generator of Section~\ref{sec:localization}.

Let $M=\Q(\zeta_7)$. For a rational conductor $q\equiv1\pmod7$, choose a generator $\pi_q$ of norm $q$ for the selected degree one prime. The source gives exact coefficient vectors in $1,\zeta_7,\ldots,\zeta_7^5$ for $q=197,337,547,673,743,883,1933$. For example,
\[
\pi_{197}=\zeta_7+2\zeta_7^2+2\zeta_7^4,
\]
and
\[
\pi_{337}=-1+\zeta_7-\zeta_7^2+3\zeta_7^3-\zeta_7^4+\zeta_7^5.
\]
The norm equality identifies one prime ideal. It does not leave an unspecified product of several primes.

At the chosen $q$ place, write $\pi_q=q u_q$ with $u_q$ a unit. The same norm identity gives
\begin{equation}\label{eq:selfunit}
\ell_q(u_q)=-\sum_{b=2}^6\ell_q(\pi_q(z_q^b)).
\end{equation}
Principal reciprocity on $\pi_q$ therefore supplies the own conductor column
\begin{equation}\label{eq:selfcolumn}
D_{q,q}=\sum_{b=2}^6(1-b)\ell_q(\pi_q(z_q^b)).
\end{equation}
For $r\ne q$, the column is
\begin{equation}\label{eq:othercolumn}
D_{q,r}=-\sum_{b\in\mathrm{reps}_r}b\ell_r(\pi_q(z_r^b)).
\end{equation}
The degree two representative convention is the one proved in \eqref{eq:orbitreciprocity}. These formulas use the same arithmetic Frobenius principal idele convention as the ordinary localization rows.

The direct independent check lifts the chosen root of $\Phi_7$ modulo $q^2$ by Hensel, then computes $\pi_q(\widetilde z_q)/q$ modulo $q$. It verifies the self unit rather than assigning its missing value from the norm relation. The norm calculation and Hensel calculation share the input generator but use different arithmetic routes.

Because each support prime is inert in $K/\Q$, restriction at a rational $q\equiv1\pmod7$ multiplies the compatible rational Frobenius value by two:
\begin{equation}\label{eq:betaF}
b_q=2\sum_rD_{q,r}s_r.
\end{equation}
For $q\equiv-1\pmod7$, the compatible squared rational Frobenius has beta value zero, from the cocycle factor $1+\chi(q)^{-1}=0$. This does not say beta is unramified there.

The ordinary inertia generator differs from the frozen twisted residue convention. Its beta value is
\[
c_q=s_q\ell_q(\operatorname{Nm}_{\F_{q^2}/\F_q}g_q)
\quad(q\equiv1\pmod7),
\]
and $c_q=s_q\ell_q(g_q)$ for $q\equiv-1\pmod7$. The conversion factors are
\begin{center}
\begin{tabular}{rrrrr}\toprule
$197:5$&$307:2$&$337:5$&$547:4$&$673:1$\\
$743:1$&$797:4$&$883:3$&$1483:1$&$1933:1$\\\bottomrule
\end{tabular}
\end{center}
Setting $c_q=s_q$ without this conversion mixes two local generators and changes the problem.

Finally, the expression $c_qz(F_q)-b_qz(I_q)$ is invariant under a compatible Frobenius lift change. Both $b_q$ and $z(F_q)$ acquire their respective inertia multiples, and the extra terms cancel. This invariance does not permit omitting the self conductor unit or changing only one normalization.
